\documentclass[12pt]{beamer}
% \usetheme{umbc4}
\usetheme{moloch}

\usepackage{amsmath,amscd}
\usepackage{fontspec}
\usepackage{unicode-math}
\usepackage{pifont}
\usepackage{xcolor}
\usepackage[dvipsnames]{xcolor}
\usepackage{listings}
\usepackage{minted}
\usepackage{tikz}
\usepackage{twemojis}
\usepackage{nicematrix}

\usetikzlibrary{fit}

\setmonofont{DejaVu Sans Mono}

% aet the title bar to black background with white text
\setbeamercolor{frametitle}{bg=black, fg=white}

% ensure other structural elements (bullets, etc.) are visible
\setbeamercolor{structure}{fg=white}

\setmainfont{Libertinus Serif}
\setsansfont{Libertinus Sans}
\setmathfont{STIX Two Math}

% define colors for syntax highlighting
\definecolor{codegreen}{rgb}{0,0.6,0}
\definecolor{codegray}{rgb}{0.5,0.5,0.5}
\definecolor{codecoffee}{rgb}{0.75,0.6,0.3}
\definecolor{question}{rgb}{1.0,0.75,0.1}
\definecolor{lightblue}{rgb}{0.55,0.75,0.9}
\definecolor{yellow}{rgb}{1.0, 1.0, 0.0}
\definecolor{codecomment}{rgb}{0.97, 0.74, 0.4}
\definecolor{codepy}{rgb}{0.9, 0.85, 0.85} %{#F8BC65}
\definecolor{funcColor}{rgb}{0.47, 0.6, 0.8}
\definecolor{midgray}{rgb}{0.4, 0.4, 0.3}
\definecolor{indexwhite}{rgb}{0.9, 0.9, 0.9}

\setbeamercolor{background canvas}{bg=black}
\setbeamercolor{normal text}{fg=white}


\newcommand{\FrameTitle}[2]{\frametitle{#1 \hfill \small\textnormal{\textcolor{midgray}{#2}}}}
\newcommand{\snl}{\vspace*{8pt}}
\newcommand{\nl}{\vspace*{1em}\\}
\newcommand{\vsp}{\vspace*{1em}}

\newcommand{\Matrix}[1]{\textbf{#1}}
\newcommand{\Vector}[1]{\textbf{#1}}
\newcommand{\Scalar}[1]{\textnormal{#1}}
\newcommand{\Quiz}[1]{{\Large\textbf{\color{question}{#1}}}}

\begin{document}
\title{\center{Applied Linear Algebra \\\vsp Linear Functions \\}}

\author{\center{Devi Prasad M \\ Manipal School of Information Sciences \\ Manipal \\}}
\date{\center{August 2026}}
\maketitle

\begin{frame}\FrameTitle{Linear Functions}{Unit 1}
    We write $f: \mathbb{R}^n \rightarrow \mathbb{R}$ to mean that the
    function $f$ maps real $n$-vectors to real numbers.

    \vsp

    Function $f$ is a scalar-valued function of $n$-vectors.

    \vsp
    \textbf{Example}\snl
    $$
        f(\Vector{x}:\mathbb{R}^n) = \Vector{1}_n^{\mathrm{T}} \, \Vector{x}
    $$
\end{frame}


\begin{frame}\FrameTitle{Linear Functions}{Unit 1}
    Consider the vector inner product function
    \[
        f(\Vector{x}) = \Vector{a}^{\textrm{T}} \Vector{x} = a_1x_1 + \cdots + a_nx_n
    \]
    where $\Vector{a}$ is a fixed $n$-vector, and $\Vector{x}$ is a variable $n$-vector.

    \snl The function $f$ is called a \emph{linear function} of $\Vector{x}$.

    \snl The vector $\Vector{a}$ is called the \emph{coefficient vector} of the linear function $f$.

    \snl $\Vector{a}$ is also called the \emph{weight vector} of the linear function $f$.
\end{frame}

\begin{frame}\FrameTitle{Linear Functions}{Unit 1}
    A linear function satisfies the following two properties:
    \begin{itemize}
        \item \emph{Additivity}. $f(\Vector{x} + \Vector{y}) = f(\Vector{x}) + f(\Vector{y})$ for all $n$-vectors $\Vector{x}$ and $\Vector{y}$.
        \item \emph{Homogeneity}. $f(\alpha \Vector{x}) = \alpha\, f(\Vector{x})$ for all $n$-vectors $\Vector{x}$ and all scalars $\alpha$.
    \end{itemize}

    \vsp The vector inner product function
    \[
        f(\Vector{x}) = \Vector{a}^{\textrm{T}} \Vector{x} = a_1x_1 + \cdots + a_nx_n, \ \ \  \Vector{a}: \mathbb{R}^n, \ \ \Vector{x}: \mathbb{R}^n
    \]
    is linear.
\end{frame}

\begin{frame}\FrameTitle{Superposition and Linearity}{Unit 1}
    \begin{definition}
        A function $f$ is \emph{linear} if it satisfies the \emph{superposition principle}:
        \[
            f(\alpha \, \Vector{x} + \beta \, \Vector{y}) = \alpha \, f(\Vector{x}) + \beta \,f(\Vector{y})
        \]
        for all $n$-vectors $\Vector{x}$ and $\Vector{y}$, and all scalars $\alpha$ and $\beta$.
    \end{definition}

    \vsp
    Given the inner product function $f(\Vector{x}) = \Vector{a}^{\textrm{T}} \Vector{x}$, we can verify that it is linear:
    \[
        \begin{array}{rll}
        f(\alpha \, \Vector{x} + \beta \, \Vector{y}) & = & \Vector{a}^{\textrm{T}} (\alpha \, \Vector{x} + \beta \, \Vector{y}) \\
                                      & = & \Vector{a}^{\textrm{T}} (\alpha \, \Vector{x}) + \Vector{a}^{\textrm{T}} (\beta \, \Vector{y}) \\
                                      & = & \alpha (\Vector{a}^{\textrm{T}} \, \Vector{x}) + \beta (\Vector{a}^{\textrm{T}} \, \Vector{y}) \\
                                      & = & \alpha f(\Vector{x}) + \beta f(\Vector{y})
        \end{array}
    \]
\end{frame}


\begin{frame}\FrameTitle{Linear Functions}{Unit 1}
\end{frame}


\begin{frame}\FrameTitle{Regression model}{Unit 1}
    An $n$-vector \Vector{x} is a \emph{feature vector} of $n$ features.

    A \emph{regression model} is the affine function of \Vector{x} given by

    \[
        \widehat{\Vector{y}} = \Vector{x}^{\textrm{T}} \beta     + \Vector{v}
    \]

    where $\beta$ is an $n$-vector and \Vector{v} is an $n$-vector.\vsp

    The entries of \Vector{x} are called the \emph{regressors}.\vsp

    $\widehat{\Vector{y}}$ is called the \emph{prediction}.\vsp
\end{frame}

\begin{frame}\FrameTitle{Regression model}{Unit 1}
\end{frame}


\begin{frame}\FrameTitle{Applied Linear Regression}{Unit 1}

Prediction of Intensive Care Unit Length of Stay in the MIMIC-IV Dataset
https://www.mdpi.com/2076-3417/13/12/6930

"Four different classification and regression models for predicting LOS in the ICU are
utilized and their performances are compared. The first is the common logistic/linear
regression model (for classification/regression problems), which incorporates a linear
combination of all features in modeling. ..."

features:
body temperature,
chloride,
creatinine,
heart rate,
mean corpuscular volume (MCV),
O2 saturation, and
respiratory rate
\end{frame}

\end{document}
